%!tex root: ./prereq.tex
\documentclass[a4paper,11pt]{article}
\usepackage{akshaynotes}

% Book's title and subtitle
\title{\Huge \textbf{A primer to foundations of algebraic geometry} }

% Author
\author{
    \textsc{Akshay Popat}
}

\begin{document}

    \maketitle
    
    \tableofcontents

    \section{Motivation}
    \subsection{Rings}
    Algebraic geometry studies geometric spaces defined by polynomial equations, using algebra to understand their geometry. For example, consider the loci of points in $\mathbb{R}^2$ satisfying the following equations:
    \begin{align*}
      x^2 + y^2 = 1, && y^2 = x^3, && xy = 0.
    \end{align*}
    These describe a circle, a cusp, and the union of the two coordinate axes, respectively. Although all three objects are defined by polynomial equations, they have rather different geometry. The circle is smooth at every point, whereas the other two have singularities at the origin: points where the geometry fails to look locally like a smooth curve. Moreover, the singularities themselves are different. The cusp has a single branch that becomes increasingly steep near the origin, whereas the two coordinate axes cross there. We would like to understand such differences precisely, and to find algebraic ways of detecting and classifying them.

    A useful perspective, not just in algebraic geometry but throughout mathematics, is that we can learn about an object by studying the maps into and out of it. In geometry, one natural thing to study is the collection of functions defined on a space. For the three examples above, we might consider polynomial functions in $x$ and $y$, and ask what algebraic structure these functions form.

    Polynomials in two variables with real coefficients form a \textit{ring}, denoted by $\mathbb{R}[X,Y]$. Addition and multiplication of polynomials make this a ring, and evaluating a polynomial at a point gives a function on $\mathbb{R}^2$. However, when we restrict our attention to the circle, the polynomial $X^2 + Y^2 - 1$ evaluates to \(0\) everywhere on the circle. We encode this algebraically, by passing from $\mathbb{R}[X,Y]$ to the \textit{quotient} \(\mathbb{R}[X,Y]/(X^2 + Y^2 - 1)\). Here, $(X^2 + Y^2 - 1)$ denotes the \textit{ideal} generated by $X^2 + Y^2 - 1$, which is the set of all polynomial multiples of this polynomial. In the quotient, any two polynomials whose difference is a multiple of $X^2 + Y^2 - 1$ represent the same element. Thus the equation defining the circle becomes a relation in the ring itself.

    More generally, if a space is defined by several polynomial equations, we can collect those equations into an ideal and form the corresponding quotient of a polynomial ring. This gives us an algebraic object that records the equations defining the space and the relations they impose on polynomial functions. Such a quotient is called a \textit{coordinate ring} in the usual setting of affine algebraic geometry. Polynomial rings, ideals, and quotient rings will therefore be among the main algebraic objects of interest to us.

    There is a deeper reason to study these rings. It turns out that algebraic properties of a coordinate ring can reveal geometric properties of the space it describes. For instance, \textit{prime} ideals will allow us to construct the points of a more general kind of geometric space, called the \textit{spectrum} of a ring. We will eventually see that this construction leads to the notion of a \textit{scheme}, which generalises the familiar spaces defined by polynomial equations.

    \subsection{Topology}
    A geometric space is not merely an arbitrary collection of points. We also want to understand how those points fit together. A \textit{topology} specifies which subsets are open, and thereby allows us to define and study continuous functions. 

    For example, suppose that \(X\) is a topological space. We can consider the continuous real-valued functions on \(X\), denoted by \(C(X)\). Pointwise addition and multiplication of such functions makes \(C(X)\) a ring. Furthermore, if \(F: X \to Y\) is a continuous map between topological spaces, then any continuous function \(f: Y \to \mathbb{R}\) induces a continuous function \(f \circ F: X \to \mathbb{R}\). Thus a map of spaces naturally gives a map between their rings of functions \(F^*: C(Y) \to C(X)\).

    In algebraic geometry, however, we are usually interested not in all continuous functions, but in functions that are compatible with the algebraic structure of the space. Polynomial functions are the simplest example. As we move to more general spaces, we will need to specify which functions are available on different parts of a space, not just on the whole space. This leads to the idea of a \textit{sheaf}. Roughly speaking, a sheaf keeps track of the functions defined on each open subset of a space, together with the maps that restrict functions to smaller open subsets. It formalises the principle that compatible local descriptions can be combined to give a global description. This is important because geometric objects are often best understood locally, while we ultimately want to study the space as a whole. The spaces studied in algebraic geometry will therefore carry two kinds of information: a topology describing how their points fit together, and a sheaf of rings describing the functions available locally. A scheme is a space equipped with precisely this kind of algebraic structure, subject to additional conditions. The aim of this reading project is to develop the ideas needed to understand this construction and why it is useful.
    
    \section{Rings}
    \subsection{Definitions and examples}
    \begin{definition}
      A \emph{ring} is a tuple \((R, +, 0, \cdot, 1)\) where:
      \begin{enumerate}[label=(\roman*)]
        \item \((R, +, 0)\) is an abelian group;
        \item \(\cdot\) is an associative binary operation on \(R\) such that it distributes over addition from both the left and the right and \(1\) is the identity.
      \end{enumerate}
    \end{definition}

    In these notes, all rings have \(0 \ne 1\). In particular, there is no ring with just a single element.

    The operation \(+\) is called the \emph{addition} in \(R\) and the operation \(\cdot\) is called \emph{multiplication}. The symbol \(\cdot\) is often omitted, with juxtaposition denoting multiplication.

    \begin{definition}
      A ring is \emph{commutative} if the multiplication is commutative.
    \end{definition}

    \begin{example}
      The integers \(\mathbb{Z}\) form a commutative ring.
    \end{example}

    \begin{example}
      In linear algebra we have seen fields such as \(\mathbb{R}\) and \(\mathbb{C}\). These are examples of commutative rings. Later we will see a definition of fields as rings satisfying certain properties.
    \end{example}

    \begin{example}
      Let \(R\) be a ring. Define the polynomial ring in one variable \(X\) with coefficients in \(R\)
      \[
        R[X] = \left\{ a_n X^n + \cdots + a_1 X + a_0 : a_0, \cdots, a_n \in R \right\}.
      \]
      \(R[X]\) is a ring.
    \end{example}

    \begin{example}
      We can iterate the above construction to make a polynomial ring in multiple variables. For example, a polynomial ring in three variables
      \[
        R[X_1, X_2, X_3] = R[X_1][X_2][X_3].
      \]
    \end{example}

    \begin{example}
      Matrices with coefficients in \(\mathbb{R}\), denoted \(M_n(\mathbb{R})\), form a ring. For \(n \ge 2\) this ring is non-commutative.
    \end{example}

    From now on, we will only work with commutative rings, and all rings are assumed to be commutative unless otherwise specified.

    \begin{definition}
      Let \(R\) be a ring. An element \(u \in R\) is a \emph{unit} if there exists \(v \in R\) such that \(uv = vu = 1\). The set of units is denoted \(R^\times\).
    \end{definition}

    \begin{exercise}
      Show that the set of units is a group under the ring multiplication.
    \end{exercise}

    \begin{definition}
      A ring \(R\) is a \emph{field} if \(R^\times = R \setminus \{0\}\).
    \end{definition}

    \begin{exercise}
      Show that integers modulo \(3\) form a ring, and show that this ring is in fact a field.
    \end{exercise}

    \begin{definition}
      Let \(R\) be a ring. An element \(r \in R \setminus \{0\}\) is a \emph{zero-divisor} if there exists an element \(s \in R \setminus \{0\}\) such that \(rs = 0\).
    \end{definition}

    \begin{definition}
      An \emph{integral domain} is a ring which has no zero-divisors.
    \end{definition}

    In other words an integral domain is a ring where if \(xy = 0\) then you can conclude that either \(x = 0\) or \(y = 0\).

    \begin{exercise}
      Show that all fields are integral domains.
    \end{exercise}

    \begin{exercise}
      Show that integers modulo \(6\) form a ring. Show that this ring is not an integral domain.
    \end{exercise}

    \begin{definition}
      Let \(R\) be a ring. An element of \(R\) is \emph{irreducible} if it is non-zero, not a unit, and cannot be written as a product of two non-units.
    \end{definition}

    \begin{definition}
      An \emph{unique factorisation domain (UFD)} is an integral domain such that every non-zero element admits a unique factorisation into irreducible elements up to reordering and multiplication by units.
    \end{definition}

    \begin{example}
      \(\mathbb{Z}\) is a unique factorisation domain.
    \end{example}

    \begin{example}
      All fields are unique factorisation domains.
    \end{example}

    \begin{exercise}
      Show that \(\mathbb{Z}[\sqrt{-5}] = \{a + b\sqrt{-5} : a,b \in \mathbb{Z}\}\) is not a unique factorisation domain. (Hint: \(6 = (1 + \sqrt{-5})(1 - \sqrt{-5})\).)
    \end{exercise}

    We mention the following without proof:

    \begin{theorem}[Gauss' lemma]
      If \(R\) is a unique factorisation domain, then \(R[X]\) is a unique factorisation domain. \qed
    \end{theorem}

    \begin{example}
      \(\mathbb{C}\) is a unique factorisation domain, so by induction and Gauss' lemma, \(\mathbb{C}[X_1, \dots, X_n]\) is a unique factorisation domain for all \(n\).
    \end{example}

    \begin{definition}
      A ring \(R\) is a \emph{Euclidean domain} if it is an integral domain and there exists a function \(\phi: R \setminus \{0\} \to \mathbb{N}\) such that: for all \(x \in R\) and \(y \in R \setminus \{0\}\), there exist \(q,r \in R\) such that \(x = qy + r\) and either \(r = 0\) or \(\phi(r) < \phi(y)\).
    \end{definition}

    \begin{exercise}
      Show that \(\mathbb{Z}\) is a Euclidean domain.
    \end{exercise}

    \begin{exercise}
      Show that all fields are Euclidean domains.
    \end{exercise}

    \begin{exercise}
      Let \(k\) be a field. Show that \(k[X]\) is a Euclidean domain.
    \end{exercise}

    \subsection{Ideals}
    \begin{definition}
      Let \(R\) be a ring. A subset \(I \subseteq R\) is an \emph{ideal} if it is an abelian group and \(ri \in I\) for all \(r \in R\) and \(i \in I\).
    \end{definition}

    \begin{example}
      For any ring \(R\), the subsets \(\{0\} \subseteq R\) and \(R \subseteq R\) are ideals.
    \end{example}

    \begin{exercise}
      Show that an ideal \(I \subseteq R\) is equal to the whole ring if and only if it contains a unit.
    \end{exercise}

    \begin{example}
      Let \(R\) be a ring and let \(x \in R\). The \emph{ideal generated by \(x\)} is the ideal
      \[
        (x) = Rx = \left\{ rx : r \in R \right\}.
      \]
    \end{example}

    By the above exercise, we have \((1) = R\).

    \begin{definition}
      An ideal \(I \subseteq R\) is \emph{proper} if \(I \ne (1)\).
    \end{definition}

    \begin{definition}
      An ideal \(I \subseteq R\) is \emph{principal} if it can be written \(I = (x)\) for some \(x \in R\).
    \end{definition}

    \begin{definition}
      A \emph{principal ideal domain (PID)} is an integral domain such that every ideal is principal.
    \end{definition}

    \begin{example}
      \(\mathbb{Z}\) is a cyclic group, so all of its subgroups are cyclic also. Therefore any ideal of \(\mathbb{Z}\) is of the form \(n \mathbb{Z}\). So \(\mathbb{Z}\) is a principal ideal domain.
    \end{example}

    \begin{exercise}
      Show that \(\mathbb{C}[X,Y]\) is not a principal ideal domain.
    \end{exercise}

    \begin{exercise}
      Show that all fields are principal ideal domains.
    \end{exercise}

    We give the following diagram without proof:
    \[
      \text{Field} \implies \text{Euclidean domain} \implies \text{PID} \implies \text{UFD} \implies \text{Integral domain}.
    \]
    Note that the only implications which need proof are that all Euclidean domains are principal ideal domains, and that all principal ideal domains are unique factorisation domains. By the diagram and the exercise, any polynomial ring in one variable with coefficients in a field is a principal ideal domain.

    The above construction can be generalised to any subset of \(R\).

    \begin{example}
      Let \(R\) be a ring and let \(S \subseteq R\). The \emph{ideal generated by \(S\)} is the ideal
      \[
        (S) = \left\{ \sum_{s \in S} r_s s : r_s \in R \text{ and all but finitely many } r_s = 0 \right\}.
      \]
    \end{example}

    \begin{definition}
      An ideal \(I \subseteq R\) is \emph{finitely generated} if it can be written \(I = (S)\) for a finite set \(S \subseteq R\).
    \end{definition}

    We can combine ideals in certain ways:

    \begin{exercise}
      Show that the intersection of ideals is an ideal.
    \end{exercise}

    \begin{exercise}
      Show that for any chain of ideals \(I_1 \subseteq I_2 \subseteq \cdots\), the union \(\bigcup_n I_n\) is an ideal.
    \end{exercise}

    \begin{example}
      Given ideals \(I,J \subseteq R\), the sum
      \[
        I + J = \left\{ i + j : i \in I, j \in J \right\}
      \]
      is an ideal.
    \end{example}

    \begin{exercise}
      Find an example of a ring \(R\) and ideals \(I\) and \(J\) in \(R\) such that
      \[
        \left\{ ij : i \in I, j \in J \right\}
      \]
      is not an ideal.
    \end{exercise}

    \begin{example}
      Given ideals \(I,J \subseteq R\), the ideal generated by products of elements in \(I\) and \(J\) is denoted
      \[
        IJ = \left(ij: i \in I, j \in J\right).
      \]
    \end{example}

    \begin{exercise}
      Let \(I\) and \(J\) be ideals. Show that \(IJ \subseteq I \cap J\).
    \end{exercise}

    \begin{definition}
      A ring is \emph{Noetherian} if all of its ideals are finitely generated.
    \end{definition}

    \begin{exercise}
      Let \(R\) be a ring such that for all chains of ideals \(I_1 \subseteq I_2 \subseteq \cdots\) in \(R\), there exists some \(N \in \mathbb{N}\) such that \(I_n = I_N\) for all \(n \ge N\). Show that \(R\) is Noetherian.
    \end{exercise}

    We mention the following without proof:

    \begin{theorem}[Hilbert basis]
      If \(R\) is a Noetherian ring, then the polynomial ring \(R[X]\) is Noetherian. \qed
    \end{theorem}

    \begin{example}
      By induction, the polynomial ring \(\mathbb{C}[X_1, \dots, X_n]\) is Noetherian.
    \end{example}

    \begin{definition}
      Let \(R\) be a ring and \(I \subseteq R\) an ideal. The \emph{quotient} of \(R\) by \(I\) is
      \[
        R/I = \left\{ r + I \; \middle| \; r \in R \right\}.
      \]
    \end{definition}

    \begin{exercise}
      Fill in the details for the ring structure on \(R/I\). In particular, prove that addition and multiplication is well-defined.
    \end{exercise}

    \begin{exercise}
      Now we know that \(\mathbb{Z}/n \mathbb{Z}\) is a ring for all \(n\). What can you say about \((\mathbb{Z}/n \mathbb{Z})^\times\)?
    \end{exercise}

    \begin{exercise}
      Let \(I \subseteq R\) be an ideal. Show that there is a bijection between ideals in \(R\) containing \(I\) and ideals of \(R/I\).
    \end{exercise}

    \begin{definition}
      An ideal \(I\) in a ring \(R\) is \emph{maximal} if \(R/I\) is a field.
    \end{definition}

    \begin{exercise}
      Let \(R\) be a ring. Show that an ideal \(I \subseteq R\) is maximal if and only if: \(I \ne (1)\) and the only ideal strictly containing \(I\) is \((1)\).
    \end{exercise}

    \begin{definition}
      An ideal \(I\) in a ring \(R\) is \emph{prime} if \(R/I\) is an integral domain.
    \end{definition}

    \begin{exercise}
      Show that all maximal ideals are prime.
    \end{exercise}

    \begin{exercise}
      Let \(R\) be a ring. Show that an ideal \(I \subseteq R\) is prime if and only if: for all \(x,y \in R\), if \(xy \in I\) then \(x \in I\) or \(y \in I\).
    \end{exercise}

    \begin{definition}
      Let \(R\) be a ring. The \emph{spectrum} of \(R\) is
      \[
        \operatorname{Spec}R = \left\{ \mathfrak{p} \subseteq R\; \middle| \; \mathfrak{p}\text{ is a prime ideal} \right\}.
      \]
    \end{definition}
    We will come back to the spectrum when we talk about the Zariski topology.

    \subsection{Homomorphisms}
    \begin{definition}
      Let \(A\) and \(B\) be rings. A \emph{ring homomorphism} \(\varphi: A \to B\) is a function from the set \(A\) to the set \(B\) such that it is a homomorphism of abelian groups, and
      \begin{enumerate}[label=(\roman*)]
        \item \(\varphi(a_1 a_2) = \varphi(a_1)\varphi(a_2)\) for all \(a_1,a_2 \in A\);
        \item \(\varphi(1_A) = 1_B\).
      \end{enumerate}
    \end{definition}

    \begin{example}
      For any ring \(R\), the identity map \(R \to R\) is a ring homomorphism.
    \end{example}

    \begin{exercise}
      Show that a composition of ring homomorphisms is a ring homomorphism.
    \end{exercise}

    \begin{definition}
      A \emph{ring isomorphism} is a ring homomorphism with an inverse which is a ring homomorphism.
    \end{definition}
    
    \begin{exercise}
      Show that a bijective ring homomorphism is a ring isomorphism.
    \end{exercise}

    When there exists an isomorphism between rings \(A\) and \(B\), we say they are \emph{isomorphic}, denoted \(A \cong B\).

    \begin{exercise}
      Let \(R\) be a ring and \(I \subseteq R\) an ideal. Show that the quotient map \(R \to R/I\) is a surjective ring homomorphism.
    \end{exercise}

    \begin{exercise}
      Write \(\operatorname{Hom}(A,B)\) for the set of ring homomorphisms from a ring \(A\) to a ring \(B\). Let \(R\) be an arbitrary ring. Show that \(\operatorname{Hom}(\mathbb{Z}, R)\) is a single element set.
    \end{exercise}

    \begin{definition}
      Let \(\varphi: A \to B\) be a ring homomorphism. The \emph{kernel} of \(\varphi\) is
      \[
        \ker \varphi = \left\{ a \in A \; \middle| \; \varphi(a) = 0 \right\}.
      \]
    \end{definition}

    \begin{exercise}
      Show that the kernel of a ring homomorphism is an ideal.
    \end{exercise}

    \begin{proposition}
      Let \(R\) be a ring and \(I \subseteq R\) an ideal. For all rings \(A\) and homomorphisms \(\varphi: R \to A\) with \(I \subseteq \ker \varphi\), there exists a unique homomorphism \(\overline{\varphi}: R/I \to A\) such that
      \[
        \begin{tikzcd}
          R \arrow[d,] \arrow[r,"\varphi"] & A \\
          R/I \arrow[ru,"\exists!\overline{\varphi}"',dashed] & {}
        \end{tikzcd}
      \]
      commutes.
    \end{proposition}
    \begin{proof}
      Given such a homomorphism \(\varphi\), in order that the diagram commute, we have no choice but to make the definition
      \begin{align*}
        \overline{\varphi}: R/I &\to A \\
        r + I &\mapsto \varphi(r).
      \end{align*}
      It is well-defined since if \(r + I = r' + I\), then \(r - r' \in I\), so \(0 = \varphi(r - r') = \varphi(r) - \varphi(r')\). It can be checked easily that it is a homomorphism.
    \end{proof}

    \begin{exercise}
      Show that a homomorphism is injective if and only if its kernel is the zero ideal.
    \end{exercise}

    \begin{theorem}[First isomorphism theorem for rings]
      Let \(\varphi: R \to A\) be a surjective ring homomorphism. Then the induced map \(\overline{\varphi}: R/\ker \varphi \to A\) is an isomorphism.
    \end{theorem}
    \begin{proof}
      Exercise.
    \end{proof}

    \begin{example}
      \(\mathbb{C}[X]/(X) \cong \mathbb{C}\).
    \end{example}

    \begin{example}
      \(\mathbb{C}[X,Y]/(X - Y) \cong \mathbb{C}[X]\).
    \end{example}

    \begin{example}
      \(\mathbb{C} \cong \mathbb{R}[X]/(X^2 + 1)\).
    \end{example}

    \begin{example}
      \emph{Dual numbers} are numbers of the form \(a + b\epsilon\), where \(a,b \in \mathbb{R}\) and \(\epsilon^2 = 0\). Dual numbers form a ring isomorphic to \(\mathbb{R}[X]/(X^2)\).
    \end{example}
    
    \begin{example}
      Let \(a,b \in \mathbb{C}\). Evaluation at \(a\) and \(b\) defines a homomorphism
      \begin{align*}
        \operatorname{ev}_{(a,b)}: \mathbb{C}[X,Y] &\to \mathbb{C} \\
        f(X, Y) &\mapsto f(a,b).
      \end{align*}
      Clearly \((X - a, Y - b)\) is an ideal contained in the kernel. Since \(\mathbb{C}[X,Y]/(X - a, Y - b) \cong \mathbb{C}\) is a field, the ideal \((X - a, Y - b)\) is maximal. Clearly \(1\) is not in the kernel, so the kernel is not the whole ring. So the kernel is equal to \((X - a, Y - b)\).
    \end{example}

    \begin{exercise}
      Give an example of a polynomial ring in one variable over a field and a non-zero polynomial in that ring which evaluates to \(0\) for every element of the field. Argue that this cannot happen if the field is infinite.
    \end{exercise}

    \section{Modules}
    \subsection{Definitions and examples}
    \begin{definition}
      Let \(R\) be a ring. An \emph{\(R\)-module} is an abelian group \(M\) along with a binary operation \(\cdot: R \times M \to M\) such that:
      \begin{enumerate}[label=(\roman*)]
        \item \(r_1 \cdot (r_2 \cdot m) = (r_1r_2) \cdot m\) for all \(r_1,r_2 \in R\) and \(m \in M\);
        \item \((r_1 + r_2) \cdot m = r_1 \cdot m + r_2 \cdot m\) for all \(r_1,r_2 \in R\) and \(m \in M\);
        \item \(r \cdot (m_1 + m_2) = r \cdot m_1 + r \cdot m_2\) for all \(r \in R\) and \(m_1,m_2 \in M\);
        \item \(1 \cdot m = m\) for all \(m \in M\).
      \end{enumerate}
    \end{definition}
    The symbol \(\cdot\) is often omitted.

    \begin{exercise}
      Show that \(0_R \cdot m = 0_M\) for all \(m \in M\).
    \end{exercise}

    \begin{exercise}
      Show that \(r \cdot 0_M = 0_M\) for all \(r \in R\).
    \end{exercise}

    \begin{example}
      Every ring \(R\) is an \(R\)-module by multiplication.
    \end{example}

    \begin{example}
      Let \(R\) be a ring. Every ideal \(I \subseteq R\) is an \(R\)-module by multiplication in \(R\).
    \end{example}

    \begin{example}
      Let \(k\) be a field. A \(k\)-module is exactly a \(k\)-vector space.
    \end{example}

    \begin{example}
      A \(\mathbb{Z}\)-module is exactly an abelian group.
    \end{example}
    Think about why this is.

    \begin{example}
      For every ring \(R\) the trivial abelian group \(\{0\}\) is an \(R\)-module. This is called the \emph{zero module}, denoted \(0\).
    \end{example}

    \begin{example}
      Let \(k\) be a field. A polynomial ring \(k[X_1, \dots, X_n]\) is a \(k\)-module.
    \end{example}

    \begin{example}
      More generally, if \(R\) is a ring and \(S \le R\) is a subring (a ring inside \(R\)), then \(R\) is an \(S\)-module by multiplication in \(R\).
    \end{example}

    \begin{exercise}
      Let \(R\) be an integral domain and \(M\) an \(R\)-module. An element \(m \in M\) is called a \emph{torsion element} if there exists \(r \in R \setminus \{0\}\) such that \(rm = 0\). We say \(M\) is \emph{torsion} if every element is a torsion element. We say \(M\) is \emph{torsion free} if \(0 \in M\) is the only torsion element. Give an example of an \(R\)-module which is torsion free. Give an example of an \(R\)-module which is torsion.
    \end{exercise}

    \begin{exercise}
      Show that the abelian group \(\mathbb{Q}/\mathbb{Z}\) is a torsion \(\mathbb{Z}\)-module.
    \end{exercise}

    \begin{example}
      Fix a linear map \(T: \mathbb{C}^n \to \mathbb{C}^n\). Then we can give \(\mathbb{C}^n\) the structure of a \(\mathbb{C}[X]\)-module, where multiplication by \(X\) applies the linear map \(T\).
    \end{example}

    \begin{exercise}
      Let \(R\) be a ring and \(I \subseteq R\) an ideal. Show that the quotient ring \(R/I\) has an \(R\)-module structure.
    \end{exercise}

    \begin{definition}
      An \(R\)-module \(M\) is \emph{generated} by a set \(S \subseteq M\) if every \(m \in M\) can be written as an \(R\)-linear combination of finitely many elements of \(S\). That is, for all \(m \in M\), we can write
      \[
        m = \sum_{s \in S} r_s s
      \]
      such that all but finitely many \(r_s\) are zero.

      The set \(S\) is called a \emph{generating set} for \(M\).
    \end{definition}

    Every module has a generating set by setting \(S = M\). This is an uninteresting and unhelpful observation. We care when our module can be generated by a small generating set.

    \begin{example}
      \(\mathbb{Z}\) is generated as a \(\mathbb{Z}\)-module by the set \(\{1\}\). More generally, a ring \(R\) is generated as an \(R\)-module by the set \(\{1\}\).
    \end{example}

    \begin{definition}
      A module is \emph{finitely generated} if it has a finite generating set.
    \end{definition}

    \begin{example}
      The polynomial ring \(\mathbb{C}[X]\) is not finitely generated as a \(\mathbb{C}\)-module.
    \end{example}

    \begin{definition}
      A \emph{basis} for an \(R\)-module \(M\) is a generating set \(S \subseteq M\) which is \(R\)-linearly independent. That is, if any finite sum \(\sum_s r_s s\), where \(s \in S\) and \(r_s \in R\), is equal to zero, then each \(r_s = 0\).

      If an \(R\)-module has a basis, it is called a \emph{free \(R\)-module}.
    \end{definition}

    \begin{example}
      Let \(k\) be a field. Any \(k\)-vector space is a free \(k\)-module.
    \end{example}

    \begin{example}
      \(\mathbb{Z}/2 \mathbb{Z}\) does not have a \(\mathbb{Z}\)-basis. Therefore, it is not a free \(\mathbb{Z}\)-module.
    \end{example}

    \subsection{Submodules}
    \begin{definition}
      Let \(M\) be an \(R\)-module. A \emph{submodule} of \(M\) is a subgroup \(N \le M\) such that \(rn \in N\) for all \(r \in R\) and \(n \in N\).
    \end{definition}
    The submodule \(N\) inherits the module structure from \(M\).

    \begin{definition}
      Let \(M\) be an \(R\)-module and \(N \le M\) a submodule. The \emph{quotient module} is the abelian group \(M/N\).
    \end{definition}

    \begin{exercise}
      Show that \(M/N\) in the above definition admits an \(R\)-module structure. In particular, show that the multiplication by elements of \(R\) is well-defined.
    \end{exercise}

    \begin{example}
      We saw that \(R\) is an \(R\)-module by multiplication. An ideal \(I \subseteq R\) is a submodule of \(R\). Then the quotient \(R/I\) as an \(R\)-module is the quotient module.
    \end{example}

    \subsection{Homomorphisms}
    \begin{definition}
      Let \(A\) and \(B\) be \(R\)-modules. An \emph{\(R\)-module homomorphism} \(f: A \to B\) is a homomorphism of abelian groups such that \(f(ra) = rf(a)\) for all \(a \in A\) and \(r \in R\).
    \end{definition}

    \begin{example}
      Let \(M\) be an \(R\)-module. The identity group homomorphism \(M \to M\) is an \(R\)-module homomorphism.
    \end{example}

    \begin{exercise}
      Show that a composition of \(R\)-module homomorphisms is an \(R\)-module homomorphism.
    \end{exercise}

    \begin{definition}
      An \emph{\(R\)-module isomorphism} is an \(R\)-module homomorphism with an inverse which is an \(R\)-module homomorphism.
    \end{definition}

    \begin{exercise}
      Show that a bijective \(R\)-module homomorphism is an \(R\)-module isomorphism.
    \end{exercise}

    If there exists an \(R\)-module isomorphism between \(R\)-modules \(A\) and \(B\), we say they are \emph{isomorphic}, denoted \(A \cong B\).

    \begin{exercise}
      Let \(M\) be an \(R\)-module and \(N \le M\) a submodule. Show that the quotient map \(M \to M/N\) is an \(R\)-module homomorphism.
    \end{exercise}

    \begin{example}
      Let \(A\) and \(B\) be \(R\)-modules. There is a unique \(R\)-module homomorphism \(A \to 0\) and a unique \(R\)-module homomorphism \(0 \to B\). Their composition \(A \to 0 \to B\) is an \(R\)-module homomorphism sending each element \(a \in A\) to \(0_B \in B\). This is called the \emph{zero map}.
    \end{example}

    \begin{example}
      Let \(k\) be a field, and let \(A\) and \(B\) be \(k\)-modules (so vector spaces over \(k\)). Then a module homomorphism \(A \to B\) is the same thing as a linear map.
    \end{example}

    \begin{definition}
      The \emph{kernel} of an \(R\)-module homomorphism \(f: A \to B\) is
      \[
        \ker f = \left\{ a \in A \; \middle| \; f(a) = 0 \right\}.
      \]
    \end{definition}

    \begin{exercise}
      Show that the kernel of an \(R\)-module homomorphism is a submodule.
    \end{exercise}

    \begin{proposition}
      Let \(M\) be an \(R\)-module and \(N \le M\) a submodule. For all \(R\)-modules \(A\) and \(R\)-module homomorphisms \(f: M \to A\), if \(N \subseteq \ker f\), then there exists a unique \(R\)-module homomorphism \(\overline{f}: M/N \to A\) such that
      \[
        \begin{tikzcd}
          M \arrow[r,"f"] \arrow[d] & A \\
          M/N \arrow[ru,"\exists!\overline{f}"',dashed] & {}
        \end{tikzcd}
      \]
      commutes.
    \end{proposition}
    \begin{proof}
      Exercise.
    \end{proof}

    \begin{theorem}[First isomorphism theorem for modules]
      Let \(f: A \to B\) be a surjective \(R\)-module homomorphism. The induced map \(\overline{f}: A/\ker f \to B\) is an \(R\)-module isomorphism.
    \end{theorem}
    \begin{proof}
      Exercise.
    \end{proof}

    \subsection{Direct sums}
    \begin{definition}
      Let \(A\) be a (possibly infinite) set indexing a family \(\{M_\alpha\}_{\alpha \in A}\) of \(R\)-modules. The \emph{direct sum} is the set of vectors
      \[
        \bigoplus_{\alpha \in A} M_\alpha = \left\{ (m_\alpha)_{\alpha \in A} : m_\alpha \in M_\alpha, \text{ and all but finitely many } m_\alpha = 0 \right\}.
      \]
    \end{definition}

    \begin{exercise}
      Show that the direct sum of \(R\)-modules is an \(R\)-module with coordinatewise addition and multiplication by scalars.
    \end{exercise}

    \begin{example}
      If \(k\) is a field, then the direct sum of \(k\)-modules is the same as the direct sum of vector spaces.
    \end{example}

    \begin{example}
      The direct sum of abelian groups is the direct sum of \(\mathbb{Z}\)-modules.
    \end{example}

    \begin{example}
      Let \(\{M_\alpha\}\) be a family of \(R\)-modules and let \(M = \bigoplus_\alpha M_\alpha\). For all \(\alpha\), there is an injective homomorphism \(\iota_\alpha: M_\alpha \to M\) which sends an element \(m \in M_\alpha\) to the vector which has \(m\) in the position with index \(\alpha\) and \(0\) everywhere else. This homomorphism is called the \emph{inclusion} of \(M_\alpha\) into \(M\).
    \end{example}

    \begin{proposition}[Universal property of direct sums]
      Let \(\{M_\alpha\}\) be a family of \(R\)-modules. Let \(M = \bigoplus_\alpha M_\alpha\). The direct sum \(M\) and the inclusions \(\iota_\alpha: M_\alpha \to M\) are, up to unique isomorphism, the unique module and homomorphisms satisfying: for all \(R\)-modules \(M'\) and families of homomorphisms \(\iota_\alpha': M_\alpha \to M'\), there exists a unique homomorphism \(f: M \to M'\) such that \(\iota_\alpha' = f \iota_\alpha\) for all \(\alpha\).
      \[
        \begin{tikzcd}
          M_1 \arrow[r,"\iota_1"] \arrow[dr,"\iota_1'"'] & M \arrow[d,"\exists!f",dashed] & M_2 \arrow[l,"\iota_2"'] \arrow[dl,"\iota_2'"] \\
          {} & M' & {}
        \end{tikzcd}
      \]
    \end{proposition}
    \begin{proof}
      Given such a module \(M'\) and homomorphisms \(\iota_\alpha'\), in order that the diagram commute, we have no choice but to define \(f\) to satisfy
      \[
        f(\iota_\alpha(m)) = \iota_\alpha'(m)
      \]
      for all \(\alpha\). Since every element of \(M\) can be written as a finite sum \(\sum_\alpha \iota_\alpha(m_\alpha)\) in a unique way, this defines a module homomorphism. Therefore, the direct sum does satisfy this property. 

      Suppose there is another module \(N\) and homomorphisms \(\varphi_\alpha: M_\alpha \to N\) satisfying this property. We get a unique map \(g: N \to M\) such that \(\iota_\alpha = g \varphi_\alpha\) for all \(\alpha\). Since \(M\) and \(\{\iota_\alpha\}\) satisfy the property, we get a unique map \(f: M \to N\) such that \(\varphi_\alpha = f \iota_\alpha\) for all \(\alpha\). Therefore, \(\iota_\alpha = g f \iota_\alpha\) for all \(\alpha\). Since \(M\) and \(\{\iota_\alpha\}\) satisfy the property, there is a unique map \(h: M \to M\) such that \(\iota_\alpha = h \iota_\alpha\) for all \(\alpha\). By uniqueness of \(h\), we have \(h = gf = \operatorname{id}_M\). A similar argument shows \(fg = \operatorname{id}_N\). Since \(f\) and \(g\) are unique, this is a unique isomorphism.

      \[
        \begin{tikzcd}
          {} & M \arrow[d,"f"] & {} \\
          M_1 \arrow[ur,"\iota_1"] \arrow[r,"\varphi_1"'] \arrow[dr,"\iota_1"'] & N \arrow[d,"g"] & M_2 \arrow[ul,"\iota_2"'] \arrow[l,"\varphi_2"] \arrow[dl,"\iota_2"] \\
          {} & M & {}
        \end{tikzcd} \qquad
        \begin{tikzcd}
          {} & M \arrow[dd,"\exists!h",dashed] & {} \\
          M_1 \arrow[ur,"\iota_1"] \arrow[dr,"\iota_1"'] & {} & M_2 \arrow[ul,"\iota_2"'] \arrow[dl,"\iota_2"] \\
          {} & M & {}
        \end{tikzcd}
      \]
    \end{proof}

    \begin{exercise}
      Show that an \(R\)-module is free if and only if it is isomorphic to a direct sum \(\bigoplus_{i \in I} R\).
    \end{exercise}

    \section{Topology}
    \subsection{Definitions and examples}
    \begin{definition}
      A topological space is a pair \((X, \tau)\) of a set \(X\) and a collection of subsets \(\tau \subseteq 2^X\) such that:
      \begin{enumerate}[label=(\roman*)]
        \item \(\varnothing \in \tau\) and \(X \in \tau\);
        \item If \(U \in \tau\) and \(V \in \tau\), then \(U \cap V \in \tau\);
        \item If \(\{U_\alpha\} \subseteq \tau\), then \(\bigcup_\alpha U_\alpha \in \tau\).
      \end{enumerate}
      We say that \(\tau\) is the \emph{topology} of \(X\), and a set \(U \in \tau\) is called an \emph{open set}.
    \end{definition}
    So the definition implies that a finite intersection of open sets is open, and an arbitrary union of open sets is open. Often we just write \(X\) to denote the topological space, and implicitly understand that \(X\) comes with some specification of which sets are open.

    \begin{definition}
      Let \(X\) be a topological space. A set \(V \subseteq X\) is \emph{closed} if its complement \(X \setminus V\) is open.
    \end{definition}

    \begin{example}
      Let \(X\) be a set and let \(\tau = 2^X\) be its power set. Then \((X, \tau)\) is a topological space. We say that \(\tau\) is the \emph{discrete topology} on \(X\). It is a topology where every subset is both open and closed. A set which is both open and closed is said to be \emph{clopen}.
    \end{example}

    \begin{example}
      Let \(X\) be a set and \(\tau = \{\varnothing, X\}\). Then \((X, \tau)\) is a topological space. We say that \(\tau\) is the \emph{indiscrete topology} on \(X\).
    \end{example}

    \begin{exercise}
      Let \(X = \{a, b\}\) and \(\tau = \{\varnothing, \{a\}, \{a,b\}\}\). Show that \((X, \tau)\) is a topological space. This is called the \emph{Sierpinski topological space}, and \(\tau\) is called the \emph{Sierpinski topology}.
    \end{exercise}

    \begin{example}
      The real numbers \(\mathbb{R}\) are a topological space, where the open sets are exactly the sets \(U \subseteq \mathbb{R}\) such that for all \(x \in U\) there exists an \(\varepsilon > 0\) such that \((x - \varepsilon, x + \varepsilon) \subseteq U\).

      The set \([0,1) \subseteq \mathbb{R}\) is neither open nor closed.
    \end{example}

    This example can be generalised to any metric space.

    \begin{example}
      Let \((X, d)\) be a metric space. Then \(X\) is a topological space, where the open sets are exactly the sets \(U \subseteq X\) such that for all \(x \in U\) there exists an \(\varepsilon > 0\) such that \(B_{\varepsilon}(x) \subseteq U\).
    \end{example}

    Thus the notion of a topological space can be thought of as a generalisation of the notion of a metric space. Many notions which exist for metric spaces exist also for topological spaces, simply by replacing balls around a point with \textit{open neighbourhoods}.

    \begin{definition}
    Let \(X\) be a topological space. A \emph{point} is an element \(x \in X\). Given a point \(x \in X\), an \emph{open neighbourhood} of \(x\) is an open set \(U \subseteq X\) such that \(x \in U\).
    \end{definition}

    \begin{example}
      Let \(X\) be a topological space, \(x \in X\) a point, and \((x_n)\) a sequence of points in \(X\). We say \(x_n\) \emph{converges} to \(x\), denoted \(x_n \to x\), if for all open neighbourhoods \(U\) of \(x\), there exists an \(N\) such that if \(n \ge N\) then \(x_n \in U\).
    \end{example}

    \begin{example}
      Let \(X = \{a,b\}\) with the Sierpinski topology. The constant sequence \(x_n = b\) converges both to \(b\) and \(a\).
    \end{example}

    \begin{exercise}
      See if you can come up with definitions for the notions of `interior point', `isolated point', `limit point' and `boundary point' in a topological space. Look them up to see if you guessed the definitions correctly.
    \end{exercise}

    \begin{proposition}
      Let \(X\) be a topological space. A set \(U \subseteq X\) is open if and only if for all points \(x \in U\) there exists an open set \(V\) such that \(x \in V \subseteq U\).
    \end{proposition}
    \begin{proof}
      Assume \(U \subseteq X\) is open. Then given \(x \in X\), we have that \(U\) is an open set such that \(x \in U \subseteq U\). Conversely, if all \(x \in U\) have an open neighbourhood \(V_x \subseteq U\), then \(U = \bigcup_{x \in U} V_x\). Since a union of open sets is open, we conclude that \(U\) is open.
    \end{proof}

    \begin{definition}
      A topological space \(X\) is \emph{Hausdorff} if for all pairs of distinct points \(x\) and \(y\) in \(X\), there exist disjoint open neighbourhoods \(U\) and \(V\) of \(x\) and \(y\) respectively.
    \end{definition}

    \begin{exercise}
      Show that the space \(X = \{a,b\}\) with the Sierpinski topology is not Hausdorff.
    \end{exercise}

    \begin{exercise}
      Show that all metric spaces are Hausdorff.
    \end{exercise}

    \subsection{A basis for a topology}
    \begin{definition}
      Let \(X\) be a set. A \emph{basis} for a topology on \(X\) is a collection of subsets \(\mathcal{B} \subseteq 2^X\) such that:
      \begin{enumerate}[label=(\roman*)]
        \item For all \(x \in X\), there exists \(B \in \mathcal{B}\) containing \(x\);
        \item For all \(B_1\) and \(B_2\) in \(\mathcal{B}\), if \(x \in B_1 \cap B_2\), then there exists \(B \in \mathcal{B}\) such that \(x \in B \subseteq B_1 \cap B_2\).
      \end{enumerate}
    \end{definition}

    \begin{exercise}
      Let \(X\) be a metric space, and let \(\mathcal{B}\) be the collection of balls:
      \[
        \mathcal{B} = \left\{ B_r(x) : x \in X\ \text{ and }\ r \in \mathbb{R}_+ \right\}.
      \]
      Show that \(\mathcal{B}\) is a basis for a topology on \(X\).
    \end{exercise}

    \begin{definition}
      Let \(X\) be a set and let \(\mathcal{B}\) be a basis for a topology on \(X\). The \emph{topology generated by \(\mathcal{B}\)} is the collection of subsets \(U \subseteq X\) such that for all \(x \in U\) there exists \(B \in \mathcal{B}\) such that \(x \in B \subseteq U\).
    \end{definition}
    We must check that the topology generated by \(\mathcal{B}\) is actually a topology on \(X\).

    \begin{exercise}
      Let \(X\) be a set. Show that the collection of singleton subsets of \(X\) is a basis for a topology on \(X\), and show that the topology generated by this basis is exactly the discrete topology on \(X\).
    \end{exercise}

    \begin{exercise}
      Let \(X\) be a metric space and \(\mathcal{B}\) the basis of balls. Show that the topology of \(X\) is exactly the topology generated by \(\mathcal{B}\).
    \end{exercise}

    \begin{proposition}
      Let \(X\) be a set and \(\mathcal{B}\) a basis for a topology on \(X\). Let \(\tau\) be the topology generated by \(\mathcal{B}\). Then \(\tau\) is a topology on \(X\).
    \end{proposition}
    \begin{proof}
      Vacuously, \(\varnothing \in \tau\). Given \(x \in X\), we know there is some \(B \in \mathcal{B}\) containing \(x\), so \(X \in \tau\).

      Given \(U,V \in \tau\), let \(x \in U \cap V\) be arbitrary. Let \(B_U \in \mathcal{B}\) be a set such that \(x \in B_U \subseteq U\), and let \(B_V \in \mathcal{B}\) be a set such that \(x \in B_V \subseteq V\). Then \(x \in B_U \cap B_V\), so there is some \(B \in \mathcal{B}\) such that \(x \in B \subseteq B_U \cap B_V \subseteq U \cap V\). Thus \(U \cap V \in \tau\).

      Given \(\{U_\alpha\} \subseteq \tau\), let \(x \in \bigcup_\alpha U_\alpha\) be arbitrary. There is an \(\alpha\) such that \(x \in U_\alpha\), so we have \(B \in \mathcal{B}\) such that \(x \in B \subseteq U_\alpha\). So \(x \in B \subseteq \bigcup_\alpha U_\alpha\). Thus \(\bigcup_\alpha U_\alpha \in \tau\).
    \end{proof}

    \begin{exercise}
      Let \(X\) be a topological space and let \(\mathcal{B}\) be a basis for its topology (that is, the topology generated by \(\mathcal{B}\) is exactly the topology on \(X\)). Show that every open set in \(X\) is a union of sets in \(\mathcal{B}\).
    \end{exercise}

    \begin{proposition}
      Let \(X\) be a topological space and let \(\mathcal{B} \subseteq 2^X\) be a collection of open sets such that for all open sets \(U \subseteq X\) and points \(x \in U\) there exists a set \(B \in \mathcal{B}\) such that \(x \in B \subseteq U\). Then \(\mathcal{B}\) is a basis for the topology of \(X\).
    \end{proposition}
    \begin{proof}
      Exercise.
    \end{proof}

    \subsection{Continuous maps}
    \begin{definition}
      Let \(X\) and \(Y\) be topological spaces. A \emph{continuous map} \(f: X \to Y\) is a function where \(f^{-1}(U) \subseteq X\) is open for all open sets \(U \subseteq Y\).
    \end{definition}

    \begin{exercise}
      Let \(X\) and \(Y\) be topological spaces, and let \(X\) have the discrete topology. Show that all functions \(f: X \to Y\) are continuous.
    \end{exercise}

    \begin{exercise}
      Let \(X\) and \(Y\) be topological spaces, and let \(Y\) have the indiscrete topology. Show that all functions \(f: X \to Y\) are continuous.
    \end{exercise}

    \begin{exercise}
      Show that the definition of continuity for topological spaces agrees with the \(\varepsilon\)-\(\delta\) definition of continuity for metric spaces. That is, if \(X\) and \(Y\) are metric spaces, show that \(f: X \to Y\) is continuous as a map of topological spaces if and only if \(f\) is continuous under the \(\varepsilon\)-\(\delta\) definition.
    \end{exercise}

    \begin{example}
      For all topological spaces \(X\), the identity map \(\operatorname{id}_X: X \to X\) is continuous.
    \end{example}

    \begin{exercise}
      Show that a composition of continuous maps is continuous.
    \end{exercise}

    \begin{definition}
      A \emph{homeomorphism} is a continuous map of topological spaces which has an inverse which is also continuous.

      If a homeomorphism exists between topological spaces \(X\) and \(Y\), we say that they are \emph{homeomorphic}, denoted \(X \cong Y\).
    \end{definition}

    \subsection{Product topology (briefly)}
    Given two topological spaces \(X\) and \(Y\), we can give a topology to the Cartesian product \(X \times Y\). The collection of sets \(U \times V\), where \(U \subseteq X\) is open and \(V \subseteq Y\) is open, forms a basis for a topology on \(X \times Y\) called the \emph{product topology}. Equivalently, the product topology is the coarsest topology (meaning, the topology with the fewest open sets) which makes the projection maps \(\pi_X: X \times Y \to X\) and \(\pi_Y: X \times Y \to Y\) continuous.

    The product satisfies the following universal property. Given a topological space \(Z\) and maps \(\lambda_X: Z \to X\) and \(\lambda_Y: Z \to Y\), there is a unique map \(f: Z \to X \times Y\) such that
    \[
      \begin{tikzcd}
        X & X \times Y \arrow[l,"\pi_X"'] \arrow[r,"\pi_Y"] & Y \\
        {} & Z \arrow[ul,"\lambda_X"] \arrow[ur,"\lambda_Y"'] \arrow[u,"\exists!f",dashed] & {}
      \end{tikzcd}
    \]
    commutes.

    \subsection{Subspace topology (briefly)}
    If \(X\) is a topological space and \(A \subseteq X\) is a subset, there is a natural way to turn \(A\) into a topological space. We say that a subset \(U \subseteq A\) is open if we can write \(U = W \cap A\) for some open set \(W \subseteq X\). This defines a topology on \(A\), called the \emph{subspace topology}. Equivalently, the subspace topology is the coarsest topology which makes the inclusion map \(\iota: A \to X\) continuous.

    For example, the open interval \((0,1) \subseteq \mathbb{R}\) with the subspace topology is a topological space homeomorphic to \(\mathbb{R}\).

    We can construct many interesting topological spaces by considering subspaces of \(\mathbb{R}^n\). For example, the \(n\)-sphere is the topological space
    \[
      S^n = \left\{ x \in \mathbb{R}^{n+1} \; \middle| \; \|x\| = 1 \right\}.
    \]
    Its topology is the subspace topology, viewed as a subspace of \(\mathbb{R}^{n+1}\).

    \(S^1\) is a circle. Using the product topology, we can form the topological space \(S^1 \times S^1\), known as the \emph{torus}. Equivalently, we could view the torus as a subspace of \(\mathbb{R}^3\).
    
    \subsection{Quotient topology (briefly)}
    Given a topological space \(X\) and a subset \(A \subseteq X\), we can form the \emph{quotient space} \(X/A\). As a set, it is the set \(X/{\sim}\) under the equivalence relation
    \[
      x \sim y \iff (x = y\ \text{ or }\ x,y \in A).
    \]
    So it is the space obtained by collapsing all of the subspace \(A\) into a single point. The topology we give to \(X/A\) is the one such that a set \(U \subseteq X/A\) is open if and only if its preimage \(q^{-1}(U)\) is open, where the map \(q: X \to X/A\) is the quotient map. This is called the \emph{quotient topology}.

    Consider, for example, the interval \([0,1]\) as a subspace (with the subspace topology) of \(\mathbb{R}\). We can glue the two endpoints together to get a circle. That is, \(S^1 \cong [0,1]/\{0,1\}\). This is another way to think about the circle as a topological space.

    Quotient spaces satisfy a property reminiscent of one we have seen for both rings and modules. If \(X\) and \(Y\) are topological spaces, \(A \subseteq X\) is a subspace, and \(f: X \to Y\) is a continuous map such that \(f(a) = f(b)\) for all \(a,b \in A\) (that is, all of \(A\) is mapped to the same point), then there is a unique induced map \(\overline{f}: X/A \to Y\) such that
    \[
      \begin{tikzcd}
        X \arrow[r,"f"] \arrow[d] & Y \\
        X/A \arrow[ru,"\exists!\overline{f}"',dashed] & {}
      \end{tikzcd}
    \]
    commutes.

    \subsection{Compactness (briefly)}
    Compactness is an important finiteness condition on topological spaces. An \emph{open cover} for a topological space \(X\) is a collection of open sets whose union is equal to all of \(X\). A space is \emph{compact} if every open cover has a finite subcover, meaning a finite subcollection of the cover whose union is still all of \(X\). As a very rough explanation, compact spaces are necessary in many situations, as they provide a way to turn infinite things, which are potentially problematic, into finite things.

    Compactness is preserved under homeomorphisms. This means that we can differentiate between spaces using compactness. For example, the interval \([0,1]\) is compact, but the real line \(\mathbb{R}\) is not. Therefore, there is no homeomorphism between them.

    \subsection{Connectedness (briefly)}
    A topological space is \emph{disconnected} if there are two disjoint non-empty open sets whose union is the whole space. For example, the subspace \(X = \mathbb{R} \setminus (-1, 1)\), as a subspace of \(\mathbb{R}\), is disconnected. Indeed, \(((-\infty, 0) \cap X) \cup ((0, \infty) \cap X) = X\).

    A topological space is \emph{connected} if it is not disconnected. Equivalently, a topological space is connected if and only if the only clopen sets are the empty set and the whole space.

    Connectedness provides another way to differentiate spaces. A connected space cannot be homeomorphic to a disconnected space.

    A similar notion exists called \textit{path-connectedness}. A topological space \(X\) is \emph{path-connected} if for any two points \(x,y \in X\), there is a continuous map \(\gamma: [0,1] \to X\), called a \emph{path}, such that \(\gamma(0) = x\) and \(\gamma(1) = y\). Every path-connected space is connected, however the converse need not hold in general. (Look up \textit{the topologist's sine curve}.)

    \subsection{Algebraic invariants (roughly)}
      Compactness and the number of connected components are examples of \textit{invariants} of topological spaces. This means they are preserved by homeomorphisms. Invariants are useful, because they provide a way to tell spaces apart. For example, if one space is connected whereas another is not connected, then the two cannot be homeomorphic. The number of connected components is an integer valued invariant. An algebraic invariant is when we associate an algebraic object, such as a group or a ring, to a topological space, and homeomorphisms induce isomorphisms.

      What we really want is a \textit{functor} from the \textit{category} of topological spaces to an algebraic \textit{category}. Roughly, a category is a collection of objects and maps between them, satisfying some properties. Examples include: the category of sets and functions between sets, the category of groups and group homomorphisms, the category of rings and ring homomorphisms, the category of \(R\)-modules and \(R\)-module homomorphisms, the category of topological spaces and continuous maps, etc. A functor is a map of categories. It maps objects of a category \(\mathcal{C}\) to objects of a category \(\mathcal{D}\), and it maps maps in \(\mathcal{C}\) to maps in \(\mathcal{D}\), in a way which preserves composition.
      \[
        \begin{tikzcd}
          A \arrow[d,"\alpha"] \arrow[r,"\beta"] & B\\
          C \arrow[ur,"\gamma"] \arrow[r,"\delta"] & D
        \end{tikzcd}
        \implies
        \begin{tikzcd}
          F(A) \arrow[d,"F\alpha"] \arrow[r,"F\beta"] & F(B) \\
          F(C) \arrow[ur,"F\gamma"] \arrow[r,"F\delta"] & F(D)
        \end{tikzcd}
      \]
      Any functor from the category of topological spaces to an algebraic category would map homeomorphisms to isomorphisms, thus it would be an algebraic invariant.

      The \textit{fundamental group} of a topological space \(X\) at a point \(x\), denoted \(\pi_1(X,x)\) is a group consisting of homotopy classes of loops in \(X\) based at \(x\). A \textit{loop based at \(x\)} is a continuous map \(\gamma: [0,1] \to X\) such that \(\gamma(0) = \gamma(1) = x\). We can concatenate two loops \(\gamma\) and \(\gamma'\) based at \(x\) to get another loop \(\gamma * \gamma'\) based at \(x\). However, the concatenation of loops is not associative. If we have three loops \(\alpha\), \(\beta\) and \(\gamma\) based at \(x\), the concatenations \((\alpha * \beta) * \gamma\) and \(\alpha * (\beta * \gamma)\) move at different speeds at different parts of the loops, so as maps \([0,1] \to X\) they are different, although they have the same image. To rectify this, we consider loops \textit{up to homotopy}. Two loops are \textit{homotopic} if one can be continually deformed into another while leaving the basepoint \(x\) fixed. Concatenation of loops respects homotopy, in that if \(\alpha\) and \(\alpha'\) are homotopic, and \(\beta\) and \(\beta'\) are homotopic, then the concatenations \(\alpha * \beta\) and \(\alpha' * \beta'\) are homotopic. Homotopy classes of loops based at \(x\) with concatenation form a group, with the identity being the homotopy class of the constant loop \(e(t) = x\). This group is called the \textit{fundamental group}, denoted \(\pi_1(X, x)\). A continuous map \(f: X \to Y\) induces a group homomorphism \(f_*: \pi_1(X,x) \to \pi_1(Y,f(x))\), by composition
      \[
        f_*([\gamma]) = [f \circ \gamma].
      \]
      This satisfies the properties \((\operatorname{id}_X)_* = \operatorname{id}_{\pi_1(X)}\) and \((f \circ g)_* = f_* \circ g_*\). Thus \(\pi_1\) is a functor.

      \textit{Homology} is a related notion defined in a more subtle way. The \emph{standard \(n\)-simplex} is the topological space
      \[
        \Delta^n = \left\{ \left(x_0, \dots, x_n\right) \in \mathbb{R}^{n+1} \; \middle| \; x_i \ge 0\text{ for all }i\ \text{ and }\ x_0 + \cdots + x_n = 1 \right\}.
      \]
      We consider, for each integer \(n \ge 0\), continuous maps \(\sigma: \Delta^n \to X\). The free abelian group generated by these maps is called the \(n^\text{th}\) chain group, denoted \(C_n(X)\). There is a way in which for a chain \(\sigma \in C_n(X)\) we can take its boundary \(\partial \sigma\), which is a chain in \(C_{n-1}(X)\) roughly corresponding to the faces of the \(n\)-simplex. The boundary operator is a group homomorphism \(C_n(X) \to C_{n-1}(X)\). Since we can do this for any \(n\), we get a \textit{chain complex}
      \[
        \begin{tikzcd}
          C_1(X) & C_2(X) \arrow[l,"\partial"'] & \cdots \arrow[l,"\partial"'] & C_n(X) \arrow[l,"\partial"'] & C_{n+1}(X) \arrow[l,"\partial"'] & \cdots \arrow[l,"\partial"']
        \end{tikzcd}
      \]
      Taking the boundary twice in a row always results in the \(0\) chain. (Think about how the boundary of a \(2\)-simplex, which is a filled in triangle, is the three edges. Since the three edges are joined up, taking the boundary again gives the empty set.) This means that the image of the boundary homomorphism is a subgroup of the kernel of the next boundary homomorphism. Since the groups are abelian, we can take the quotient of the kernel by the image. This quotient is the \textit{Homology group}, denoted \(H_n(X)\). A \textit{homology class} is an element \([\sigma] \in H_n(X)\). This means that \(\partial \sigma = 0\), and \([\sigma] = [\sigma']\) if and only if \(\sigma - \sigma' = \partial \tau\) for some \(\tau \in C_{n+1}(X)\). Since each chain group is generated by maps \(\sigma: \Delta^n \to X\), a continuous map \(f: X \to Y\) induces, for all \(n\), a homomorphism of chain groups \(f_\sharp: C_n(X) \to C_n(Y)\) by composition, \(f_\sharp(\sigma) = f \circ \sigma\). This commutes with taking the boundary. That is, \(f_\sharp(\partial \sigma) = \partial f_\sharp(\sigma)\). Therefore, for all \(n\), the image of \(\partial: C_n(X) \to C_{n-1}(X)\) is mapped by \(f_\sharp\) into the image of \(\partial: C_n(Y) \to C_{n-1}(Y)\). Therefore, we get well-defined induced homomorphisms on the quotients \(f_*: H_n(X) \to H_n(Y)\). This satisfies \((\operatorname{id}_X)_* = \operatorname{id}_{H_n(X)}\) and \((f \circ g)_* = f_* \circ g_*\). Thus \(H_n\) is a functor, from the category of topological spaces to the category of abelian groups.

      Dual to homology is the notion of \textit{cohomology}. The \textit{\(n^\text{th}\) cochain group} is the abelian group of group homomorphisms from \(C_n(X)\) to \(\mathbb{Z}\), denoted
      \[
        C^n(X) = \operatorname{Hom}_{\mathbb{Z} }(C_n(X), \mathbb{Z}).
      \]
      The boundary homomorphisms \(\partial\) induce \textit{coboundary homomorphisms} \(\delta: C^n(X) \to C^{n+1}(X)\) by composition, \(\delta(f)(\sigma) = f(\partial \sigma)\). Thus we get a \textit{cochain complex}
      \[
        \begin{tikzcd}
          C^1(X) \arrow[r,"\delta"] & C^2(X) \arrow[r,"\delta"] & \cdots \arrow[r,"\delta"] & C^n(X) \arrow[r,"\delta"] & C^{n+1}(X) \arrow[r,"\delta"] & \cdots
        \end{tikzcd}
      \]
      Taking quotients of kernels by images yields the \textit{cohomology} groups \(H^n(X)\). This time a continuous map \(f: X \to Y\) induces homomorphisms \(f^*: H^n(Y) \to H^n(X)\). This satisfies \((\operatorname{id}_X)^* = \operatorname{id}_{H^n(X)}\) and \((f \circ g)^* = g^* \circ f^*\). We say that \(H^n\) is a \textit{contravariant functor}. In fact, there is a way to take a product of cohomology classes to get another cohomology class, called the \textit{cup product}.
      \[
        \cup: H^m(X) \times H^n(X) \to H^{m+n}(X)
      \]
      This means that if we take a direct sum of the cohomology groups
      \[
        H^*(X) = \bigoplus_n H^n(X)
      \]
      we get a \textit{graded ring}. The induced map \(f^*\) is a graded ring homomorphism \(f^*: H^*(Y) \to H^*(X)\), so \(H^*\) is a contravariant functor, from the category of topological spaces to the category of graded rings.

    \subsection{The ring of continuous functions}
    Let \(X\) be a non-empty topological space. We write \(C(X)\) to denote the set of continuous functions from \(X\) to \(\mathbb{R}\).

    \begin{exercise}
      Show that a constant function from a topological space \(X\) to \(\mathbb{R}\) is always continuous.
    \end{exercise}

    By pointwise addition and pointwise multiplication, we have that \(C(X)\) forms a ring, where the constant functions \(0\) and \(1\) are the additive and multiplicative identities, respectively.

    \begin{exercise}
      Show that the pointwise sum and product of continuous functions is continuous.
    \end{exercise}

    Let \(F: X \to Y\) be a continuous map of topological spaces. There is an induced map
    \begin{align*}
      F^*: C(Y) &\to C(X) \\
      f &\mapsto f \circ F
    \end{align*}

    \begin{exercise}
      Show that the map \(F^*: C(Y) \to C(X)\) is a ring homomorphism.
    \end{exercise}

    \begin{exercise}
      Show that \((\operatorname{id}_X)^* = \operatorname{id}_{C(X)}\), and \((F \circ G)^* = G^* \circ F^*\).
    \end{exercise}

    Thus \(C\) is a contravariant functor, from the category of non-empty topological spaces to the category of rings.

    \subsection{The Zariski topology}
    We return to the spectrum of a ring. Given a ring \(R\), recall
    \[
      \operatorname{Spec}R = \left\{ \mathfrak{p} \subseteq R \; \middle| \; \mathfrak{p}\text{ is a prime ideal} \right\}.
    \]
    The \emph{Zariski topology} is a topology on \(\operatorname{Spec}R\), where the closed sets are exactly the subsets
    \[
      V(I) = \left\{ \mathfrak{p} \in \operatorname{Spec}R \; \middle| \; I \subseteq \mathfrak{p} \right\}
    \]
    for all ideals \(I \subseteq R\).

    \begin{exercise}
      Show that this defines a topology on \(\operatorname{Spec}R\).
    \end{exercise}

    \begin{exercise}
      Show that a singleton \(\{\mathfrak{p}\}\) is closed if and only if \(\mathfrak{p}\) is maximal.
    \end{exercise}
    
    \begin{example}
      Let us think about the space \(\operatorname{Spec}\mathbb{C}[X]\). Since \(\mathbb{C}[X]\) is a principal ideal domain, its ideals are of the form \((f)\), where \(f \in \mathbb{C}[X]\) is a polynomial. Suppose \((f)\) is a prime ideal. If \(f = gh\), since \(gh \in (f)\) and \((f)\) is prime, either \(g \in (f)\) or \(h \in (f)\). Either \(g\) or \(h\) has degree greater than or equal to \(\operatorname{deg}f\). Since \(g\) and \(h\) are factors of \(f\), if \(f \ne 0\), they cannot have higher degree. Therefore, if \(f \ne 0\), one of \(g\) or \(h\) has the same degree as \(f\), while the other has degree \(0\). A degree \(0\) polynomial in \(\mathbb{C}[X]\) is a unit. Thus, either \(f = 0\) or \(f\) is irreducible.

      Since \(\mathbb{C}[X]\) is an integral domain, the ideal \((0)\) is prime. If \(f \in \mathbb{C}[X]\) is an irreducible non-zero polynomial, then \((f)\) is a prime ideal since \(\mathbb{C}[X]\) is a unique factorisation domain. Since \(\mathbb{C}\) is algebraically closed, the irreducible polynomials are non-zero scalar multiples of \(X - a\), where \(a \in \mathbb{C}\). Thus the points of \(\operatorname{Spec}\mathbb{C}[X]\) are exactly those of the form \((X - a)\), and the point \((0)\).

      The topology is strange. Every non-empty open set contains the point \((0)\). It is called a \textit{generic point}. In particular, the space is not Hausdorff. Since \(\mathbb{C}[X]\) is a principal ideal domain, the closed sets are of the form \[V(f) = \{(X - a) : (X - a)\text{ divides }f\}.\] Each \(f\) only has finitely many factors, so all closed sets are finite. All finite sets are closed, since given a finite set \(\{(X - a_1), \dots, (X - a_n)\}\) can be written
      \[
        \left\{a_1, \dots, a_n\right\} = V((X - a_1) \cdots (X - a_n)).
      \]
    \end{example}
    \nocite{*}
    \printbibliography

\end{document}
